\documentclass[a4paper,10pt]{article}
\usepackage[francais]{babel}
\usepackage{url}

\selectlanguage{francais}
\input{letter-ensiie.sty}

\selectlanguage{francais}
\enCouleur

\begin{document}


\name{St\'ephane \textsc{Genaud}}
\fonction{Directeur de l'antenne de Strasbourg de l'ENSIIE\\Professeur des Universités, ENSIIE}
\telephoneillkirch{+33 (0)3688 54542}
\email{genaud@ensiie.fr}

\begin{letter}[Recommandation de Hicham El Bilali pour une demande de naturalisation]
	  {à}{M. Le Préfet}


Monsieur le Préfet,

L'ENSIIE est une école d'ingénieur publique proposant un cursus en 3 ans en informatique.
Hicham El Bilali a intégré l'école sur le concours Centrale-Supélec après deux ans de classes préparatoires.
En tant que directeur de l'antenne de l'ENSIIE à Strasbourg, j'ai eu le plaisir de compter 
parmi mes étudiants M. Hicham El Biblali pour l'année scolaire 2012/2013. Il a ensuite
poursuivi son cursus dans le domaine de l'informatique financière sur le site d'EVry de 
l'ENSIIE.

Etant également enseignant, j'ai eu Hicham dans mes classes de cours. Je peux témoigner du
sérieux, de l'assuidité et de l'investissement sans faille d'Hicham. C'est une personne
qui s'est parfaitement intégrée dans son environnement, à la fois estudiantin et professionnel,
attesté par ses soutenances de stage (le cursus compte trois stages).

Hicham a obtenu le diplôme d'ingénieur de l'ENSIIE après un parcours scolaire sans faute.
Il a ensuite été embauché immédiatement par l'entreprise MUREX qui opère dans le domaine financier.
Hicham apporte donc des compétences informatiques dans un domaine clé de l'économie.

Pour toutes ces raisons, je soutiens au plus haut point la demande de naturalisation de M. Hicham El Bilali.

Je vous prie, M. Le Préfet, d'agréer mes salutations distinguées.

\end{letter}
\begin{flushright}
%\includegraphics[width=.2\textwidth]{signgenaud.jpg}
\end{flushright}
\end{document}
